% =====================================================================
% EXERGUE ET REMERCIEMENTS
%
% Fichier PARTAGÉ par les quatre versions du document, au même titre que les
% chapitres et que la page de garde : main.tex, main_v2.tex, main_v3.tex et
% main_ensae.tex l'appellent tous. NE PAS LE DUPLIQUER pour faire évoluer une
% version, ce serait quatre textes à maintenir.
%
% UN CHAMP RESTE À COMPLÉTER, ET IL N'APPARTIENT PAS À L'ASSISTANT : le nom de
% famille du fondateur de Nexialog Consulting. Il s'imprime en gras entre
% crochets tant qu'il n'est pas renseigné, donc un dépôt avec un nom absent est
% visible à la relecture plutôt que silencieux.
% =====================================================================

% --- Le champ à compléter --------------------------------------------
\newcommand{\fondateurnexialog}{Ali~BEHBAHANI}
% ---------------------------------------------------------------------

% =====================================================================
%  EXERGUE
%
%  La citation est la « loi de la crédibilité décroissante » de Manski, qui
%  ouvre Identification Problems in the Social Sciences. Elle est à sa place
%  ici et pas seulement par ornement : le mémoire refuse de poser la direction
%  de la contagion et publie un capital en bornes, ce qui est exactement
%  l'arbitrage qu'elle énonce. L'original anglais est :
%    « The credibility of inference decreases with the strength of the
%      assumptions maintained. »
%  La traduction est de l'auteur, et le document le dit.
% =====================================================================
\thispagestyle{empty}
\vspace*{0.30\textheight}

\begin{flushright}
\begin{minipage}{0.72\textwidth}
\itshape
La crédibilité d'une inférence décroît avec la force des hypothèses que l'on
maintient.

\vspace{10pt}
\normalfont\small\raggedleft
Charles F. Manski, \emph{Identification Problems in the Social Sciences}, 1995.\\
\textup{Traduction de l'auteur.}
\end{minipage}
\end{flushright}

\clearpage

% =====================================================================
%  REMERCIEMENTS
% =====================================================================
\chapter*{Remerciements}
\addcontentsline{toc}{chapter}{Remerciements}
% Voir la note de 20_notes_synthese : un \chapter* laisse le titre courant sur
% la marque precedente. Sans effet visible ici, la page etant unique, mais la
% regle vaut pour les quatre versions et le texte peut s'allonger.
\markboth{Remerciements}{}

Ce mémoire doit beaucoup à des relectures qui ont porté sur ce qu'il affirmait
plutôt que sur la façon dont il le calculait, et c'est de là que sont venues la
plupart de ses corrections.

Je remercie Hugo Rapior, mon maître de stage, pour un point d'avancement
hebdomadaire tenu sans exception pendant six mois. Aucune de ses demandes ne
portait sur un calcul et chacune a fait tomber un défaut : le retrait des
conclusions tirées d'un intervalle trop large, la table complète des leviers dont
est sortie la lecture de fermeture, et deux mots mal employés dont la correction a
révélé qu'une borne était annoncée à l'envers.

Je remercie Caroline Hillairet, ma tutrice académique, dont les remarques ont
toutes porté sur le traitement de l'incertitude. La plus féconde, sur le quantile
d'un objet incertain, a produit deux briques du modèle et un résultat qui la
contredit sur le mécanisme tout en lui donnant raison sur le fond. Le traitement
des processus auto-excités, qui relève de son domaine de recherche, a été testé
contre ses propres variantes de noyau plutôt que contre un modèle de paille, et
s'en est trouvé reformulé.

Je remercie \fondateurnexialog, fondateur de Nexialog Consulting, de m'avoir
accueilli dans un cabinet où un stagiaire dispose du temps long nécessaire à un
travail de cette nature, ainsi que l'ensemble des équipes du pôle Actuariat, et en
particulier la consultante qui a accepté de contredire sept affirmations du modèle
et dont la réponse a illustré une limite au lieu de la lever.

Je remercie l'ENSAE Paris et les enseignants de la filière Actuariat. La
reconnaissance qu'un paramètre puisse ne pas être \emph{identifié} plutôt que mal
estimé vient d'un de leurs cours, et c'est elle qui a changé la nature de ce
travail.

Je remercie enfin ma famille, dont le soutien n'a pas eu besoin d'être expliqué.

\clearpage
